% Editable publishing source. Rebuild without --refresh-tex to preserve edits.
\documentclass[10pt,a4paper]{article}
% Shared editable typography for all five document editions.
\usepackage{fontspec}
\setmainfont{DejaVu Serif}
\setsansfont{DejaVu Sans}
\setmonofont{DejaVu Sans Mono}[Scale=.86]
\usepackage{polyglossia}
\setdefaultlanguage{russian}
\setotherlanguage{english}
\usepackage[a4paper,margin=18mm,top=21mm,bottom=21mm,headheight=15pt]{geometry}
\usepackage{xcolor}
\definecolor{Accent}{HTML}{185A86}
\definecolor{Muted}{HTML}{526573}
\usepackage{graphicx}
\usepackage{longtable,booktabs,array,calc}
\usepackage{fvextra}
\usepackage{pdflscape}
\usepackage{titlesec}
\usepackage{fancyhdr}
\usepackage{enumitem}
\usepackage{etoolbox}
\usepackage{needspace}
\usepackage{tocloft}
\usepackage[unicode,colorlinks=true,linkcolor=Accent,urlcolor=Accent,bookmarksnumbered=false]{hyperref}
\usepackage{bookmark}
\usepackage{microtype}
\AtBeginDocument{\raggedright}
\setcounter{secnumdepth}{0}
\setcounter{tocdepth}{1}
\setlength{\parindent}{0pt}
\setlength{\parskip}{4pt plus 1pt}
\linespread{1.0}
\setlength{\emergencystretch}{3em}
\setlength{\tabcolsep}{4pt}
\renewcommand{\arraystretch}{1.16}
\AtBeginEnvironment{longtable}{\small}
% Measure a compact table or its header and first two data rows. The output
% remains a breakable longtable. Cap the reservation for unusually tall rows.
\newsavebox{\DocumentTableStart}
\newlength{\DocumentTableSpace}
\newcommand{\KeepTableStart}[2]{%
  \setbox\DocumentTableStart=\vbox{\hsize=\linewidth\small #2}%
  \setlength{\DocumentTableSpace}{\dimexpr\ht\DocumentTableStart+\dp\DocumentTableStart
    +\LTpre+\baselineskip+#1\baselineskip*4\relax}%
  \ifdim\DocumentTableSpace>.5\textheight
    \setlength{\DocumentTableSpace}{.5\textheight}%
  \fi
  \Needspace{\DocumentTableSpace}%
}
% Measure short lists once. Small blocks move together; long ones are unboxed
% into the normal page builder and remain breakable, never scaled or clipped.
\newsavebox{\DocumentShortBlock}
\newcommand{\KeepShortBlock}[1]{%
  % A heading's page-space check must not run inside the measuring box: it
  % would inspect the outer page and embed a spurious forced break in the box.
  \par\setbox\DocumentShortBlock=\vbox{\hsize=\linewidth
    \renewcommand{\Needspace}[1]{}#1\par}%
  \ifdim\dimexpr\ht\DocumentShortBlock+\dp\DocumentShortBlock\relax<.32\textheight
    \Needspace{\dimexpr\ht\DocumentShortBlock+\dp\DocumentShortBlock+\medskipamount\relax}%
  \else
    \Needspace{5\baselineskip}%
  \fi
  % Unboxing does not add normal inter-paragraph baseline glue at the leading
  % edge. Give the measured block standard vertical separation from prior text.
  \addvspace{\medskipamount}%
  \unvbox\DocumentShortBlock
}
\preto\section{\Needspace{5\baselineskip}}
\preto\subsection{\Needspace{5\baselineskip}}
\preto\subsubsection{\Needspace{4\baselineskip}}
\clubpenalty=10000
\widowpenalty=10000
\setlist{leftmargin=1.5em,itemsep=2pt,topsep=3pt}
\providecommand{\tightlist}{\setlength{\itemsep}{1pt}\setlength{\parskip}{0pt}}
\titleformat{\section}{\large\sffamily\bfseries\color{Accent}}{}{0pt}{}
\titleformat{\subsection}{\normalsize\sffamily\bfseries\color{Accent}}{}{0pt}{}
\titleformat{\subsubsection}{\normalsize\sffamily\bfseries\color{Accent}}{}{0pt}{}
\titlespacing*{\section}{0pt}{18pt}{8pt}
\titlespacing*{\subsection}{0pt}{13pt}{6pt}
\titlespacing*{\subsubsection}{0pt}{10pt}{5pt}
\pagestyle{fancy}
\fancyhf{}
\fancyfoot[C]{\small\color{Muted}\thepage}
\renewcommand{\headrulewidth}{.35pt}
\renewcommand{\footrulewidth}{0pt}
% A single compact opening page, with profile-controlled TOC depth. Reserve a
% wider page-number column for the multi-page BRD and example editions.
\setlength{\cftbeforesecskip}{5pt}
\setlength{\cftbeforesubsecskip}{2pt}
\renewcommand{\cftsecfont}{\small\sffamily\bfseries}
\renewcommand{\cftsubsecfont}{\small\sffamily}
\renewcommand{\cftsecpagefont}{\small}
\renewcommand{\cftsubsecpagefont}{\small}
\cftsetpnumwidth{2em}
\cftsetrmarg{3em}
\newcommand{\DocumentTitle}[3]{%
  \thispagestyle{plain}
  {\sffamily\bfseries\color{Accent}\fontsize{23}{29}\selectfont #1\par}
  \vspace{3mm}
  {\sffamily\small\color{Muted}#2\par}
  \ifstrempty{#3}{}{\vspace{4mm}#3\par}%
  \begingroup
  \renewcommand{\contentsname}{Содержание}%
  \tableofcontents
  \endgroup
  \vspace{5mm}
  {\color{Accent}\rule{\linewidth}{.8pt}\par}
  \clearpage
}
% Compact standalone profile: keep the title and body on the same opening page.
% It intentionally has no table of contents or forced title-page break.
\newcommand{\CompactDocumentTitle}[3]{%
  \thispagestyle{plain}
  {\sffamily\bfseries\color{Accent}\fontsize{23}{29}\selectfont #1\par}
  \vspace{3mm}
  {\sffamily\small\color{Muted}#2\par}
  \ifstrempty{#3}{}{\vspace{4mm}#3\par}%
  \vspace{5mm}
  {\color{Accent}\rule{\linewidth}{.8pt}\par}
  \vspace{4mm}
}

\hypersetup{pdftitle={Система формирования отчётности},pdfauthor={},pdfsubject={Сквозной пример}}
\fancyhead[L]{\small SolutionArchitectAI}
\fancyhead[R]{\small Сквозной пример}
\begin{document}
\DocumentTitle{Система формирования отчётности}{Сквозной пример}{Синтетические бизнес-требования, архитектурный результат, две диаграммы и полный YAML-манифест.}
\section{Исходные требования и контекст}\label{ux441ux438ux43dux442ux435ux442ux438ux447ux435ux441ux43aux438ux439-brd-ux441ux438ux441ux442ux435ux43cux430-ux444ux43eux440ux43cux438ux440ux43eux432ux430ux43dux438ux44f-ux43eux442ux447ux451ux442ux43dux43eux441ux442ux438}

Все системы, правила и числовые характеристики ниже заданы для примера. Вводные полны для проектирования данного примера; уточняющий диалог не требуется и не моделируется. {Ожидаемый результат} подготовлен по этим вводным.

\subsection{Назначение и границы}\label{ux43dux430ux437ux43dux430ux447ux435ux43dux438ux435-ux438-ux433ux440ux430ux43dux438ux446ux44b}

Аналитик просматривает продажи по подразделениям за выбранный день, скачивает сводку в CSV и передаёт её назначенному внешнему получателю. Система работает во внутреннем тестовом контуре и имеет обычную критичность; персональных данных нет.

В проект входят сервис загрузки данных и сервис формирования отчётности. Система продаж, система учёта подразделений, корпоративный брокер Apache Kafka, система корпоративного доступа, платформа PostgreSQL, система передачи отчётов и внешняя система-получатель уже существуют. Логическая база отчётности создаётся на готовой платформе. Технологический и бизнес-прокси также предоставлены корпоративным контекстом. Разработка существующих активов и физическая инфраструктура в проект не входят.

\KeepTableStart{1}{%
\begin{tabular}{@{}
  >{\raggedright\arraybackslash}p{(\columnwidth - 2\tabcolsep) * \real{0.2500}}
  >{\raggedright\arraybackslash}p{(\columnwidth - 2\tabcolsep) * \real{0.7500}}@{}}
\toprule\noalign{}
\begin{minipage}[b]{\linewidth}\raggedright
ID
\end{minipage} & \begin{minipage}[b]{\linewidth}\raggedright
Требование
\end{minipage} \\
\midrule\noalign{}
R-01 & Получать новые продажи в течение дня через корпоративный брокер Kafka. Целевая задержка появления продажи в отчётности --- не более минуты при штатной работе; при сбое или паузе по R-06 действует порядок восстановления, а минутная задержка временно не применяется \\
\end{tabular}%
}
\subsection{Требования}\label{ux442ux440ux435ux431ux43eux432ux430ux43dux438ux44f}

\begin{longtable}[]{@{}
  >{\raggedright\arraybackslash}p{(\columnwidth - 2\tabcolsep) * \real{0.2500}}
  >{\raggedright\arraybackslash}p{(\columnwidth - 2\tabcolsep) * \real{0.7500}}@{}}
\toprule\noalign{}
\begin{minipage}[b]{\linewidth}\raggedright
ID
\end{minipage} & \begin{minipage}[b]{\linewidth}\raggedright
Требование
\end{minipage} \\
\midrule\noalign{}
\endhead
\bottomrule\noalign{}
\endlastfoot
R-01 & Получать новые продажи в течение дня через корпоративный брокер Kafka. Целевая задержка появления продажи в отчётности --- не более минуты при штатной работе; при сбое или паузе по R-06 действует порядок восстановления, а минутная задержка временно не применяется \\
R-02 & Получать полный справочник подразделений каждый день в 06:00; закончить обновление до 06:15. Перед первым чтением продаж загрузить справочник \\
R-03 & Хранить продажи за последние 30 календарных дней, включая текущий день, по UTC. Загрузчик ежедневно в 00:05 удаляет продажи за пределами окна; запросы сразу ограничиваются этим окном. Повторно доставленные продажи старше окна не сохраняются. Повтор сообщения с тем же номером продажи не создаёт дубль. Дату продажи задаёт источник \\
R-04 & Аналитик выбирает день в пределах 30-дневного окна хранения и видит подразделения с продажами: наименование, количество продаж и общую сумму. При отсутствии продаж выводится пустая сводка. Запрос даты вне окна отклоняется с сообщением о допустимом диапазоне \\
R-05 & Скачивать сводку в CSV с колонками «Дата», «Подразделение», «Количество продаж», «Сумма». Файл рассчитывается на момент запроса и при пустой сводке содержит только заголовок \\
R-06 & При сбое загрузки сохранять уже полученные данные и показывать состояние загрузки и время последней успешной обработки. Не терять непрочитанные сообщения; при неизвестном подразделении приостановить обработку, показать причину и уведомить эксплуатацию. Эксплуатация исправляет справочник, загрузчик перечитывает его и продолжает с неподтверждённого сообщения. Исправление и обработка накопившихся сообщений завершаются не позднее 4 часов с начала сбоя \\
R-07 & Аналитик передаёт сводку в CSV назначенной внешней системе через корпоративную систему передачи отчётов. Пользователь видит подтверждение приёма внешним получателем или статус «Передача не подтверждена»; автоматические повторы не требуются \\
N-01 & До 5 000 продаж за день, 20 подразделений и 10 одновременных аналитиков. Целевое время ответа сводки или CSV --- не более 5 секунд \\
N-02 & Пользовательский сервис доступен в рабочие дни с 08:00 до 18:00 на уровне 99,5\%. Целевое восстановление --- до 4 часов, допустимая потеря данных --- до 24 часов \\
N-03 & Продажи и сводка содержат конфиденциальные финансовые данные без персональных сведений; справочник и состояние загрузки --- внутренние служебные данные. Каждый пользовательский запрос разрешён только сотруднику с подтверждённой ролью «Аналитик» \\
\end{longtable}

Продажи неизменяемы после передачи; исправления и возвраты в примере не рассматриваются. На первом запуске читаются новые сообщения; историческая загрузка не требуется. Разные запросы могут включать новые продажи, поэтому ранее открытая сводка и более поздний CSV могут отличаться. Все суммы заданы в одной условной валюте. Даты и расписания заданы по UTC. При восстановлении загрузчик повторно читает сообщения Kafka с позиции, соответствующей восстановленному состоянию БД, с проверкой дублей. Резервные копии хранятся 7 суток только для восстановления; 30-дневное ограничение относится к рабочему набору. До открытия восстановленной БД для запросов загрузчик удаляет из неё продажи вне текущего окна.

\KeepTableStart{1}{%
\begin{tabular}{@{}
  >{\raggedright\arraybackslash}p{(\columnwidth - 2\tabcolsep) * \real{0.2500}}
  >{\raggedright\arraybackslash}p{(\columnwidth - 2\tabcolsep) * \real{0.7500}}@{}}
\toprule\noalign{}
\begin{minipage}[b]{\linewidth}\raggedright
Набор
\end{minipage} & \begin{minipage}[b]{\linewidth}\raggedright
Содержание
\end{minipage} \\
\midrule\noalign{}
Данные о продажах & Уникальный номер продажи, дата, код подразделения и сумма \\
\end{tabular}%
}
\subsection{Данные}\label{ux434ux430ux43dux43dux44bux435}

\begin{longtable}[]{@{}
  >{\raggedright\arraybackslash}p{(\columnwidth - 2\tabcolsep) * \real{0.2500}}
  >{\raggedright\arraybackslash}p{(\columnwidth - 2\tabcolsep) * \real{0.7500}}@{}}
\toprule\noalign{}
\begin{minipage}[b]{\linewidth}\raggedright
Набор
\end{minipage} & \begin{minipage}[b]{\linewidth}\raggedright
Содержание
\end{minipage} \\
\midrule\noalign{}
\endhead
\bottomrule\noalign{}
\endlastfoot
Данные о продажах & Уникальный номер продажи, дата, код подразделения и сумма \\
Справочник подразделений & Код и наименование подразделения; коды и названия стабильны в пределах 30-дневного окна \\
Данные отчётности & Сохранённые продажи, справочник и служебное состояние загрузки \\
Сводный отчёт & Дата, наименование подразделения, количество продаж и общая сумма \\
Состояние загрузки & Состояние обработчика и время последней успешной обработки; сохраняется вместе с данными \\
Подтверждение передачи & Результат приёма внешним получателем и номер подтверждения; возвращается пользователю, отдельно не сохраняется \\
Результат проверки доступа & Разрешение на запрос и роль сотрудника; не сохраняется в БД отчётности \\
\end{longtable}

\KeepTableStart{1}{%
\begin{tabular}{@{}
  >{\raggedright\arraybackslash}p{(\columnwidth - 6\tabcolsep) * \real{0.2500}}
  >{\raggedright\arraybackslash}p{(\columnwidth - 6\tabcolsep) * \real{0.4400}}
  >{\raggedright\arraybackslash}p{(\columnwidth - 6\tabcolsep) * \real{0.2200}}
  >{\raggedright\arraybackslash}p{(\columnwidth - 6\tabcolsep) * \real{0.0900}}@{}}
\toprule\noalign{}
\begin{minipage}[b]{\linewidth}\raggedright
Система или платформа
\end{minipage} & \begin{minipage}[b]{\linewidth}\raggedright
Доступная возможность
\end{minipage} & \begin{minipage}[b]{\linewidth}\raggedright
Протокол
\end{minipage} & \begin{minipage}[b]{\linewidth}\raggedright
Порт
\end{minipage} \\
\midrule\noalign{}
Система продаж & Публикует каждую завершённую продажу в корпоративный брокер & Kafka поверх TLS & 9093 \\
\end{tabular}%
}
\subsection{Корпоративный контекст}\label{ux43aux43eux440ux43fux43eux440ux430ux442ux438ux432ux43dux44bux439-ux43aux43eux43dux442ux435ux43aux441ux442}

\begin{longtable}[]{@{}
  >{\raggedright\arraybackslash}p{(\columnwidth - 6\tabcolsep) * \real{0.2500}}
  >{\raggedright\arraybackslash}p{(\columnwidth - 6\tabcolsep) * \real{0.4400}}
  >{\raggedright\arraybackslash}p{(\columnwidth - 6\tabcolsep) * \real{0.2200}}
  >{\raggedright\arraybackslash}p{(\columnwidth - 6\tabcolsep) * \real{0.0900}}@{}}
\toprule\noalign{}
\begin{minipage}[b]{\linewidth}\raggedright
Система или платформа
\end{minipage} & \begin{minipage}[b]{\linewidth}\raggedright
Доступная возможность
\end{minipage} & \begin{minipage}[b]{\linewidth}\raggedright
Протокол
\end{minipage} & \begin{minipage}[b]{\linewidth}\raggedright
Порт
\end{minipage} \\
\midrule\noalign{}
\endhead
\bottomrule\noalign{}
\endlastfoot
Система продаж & Публикует каждую завершённую продажу в корпоративный брокер & Kafka поверх TLS & 9093 \\
Корпоративный брокер сообщений & Использует Apache Kafka; доступен издателю и загрузчику. Хранит события не менее 7 суток; повторная доставка возможна & Kafka поверх TLS & 9093 \\
Система учёта подразделений & Предоставляет полный справочник к 05:45, содержащий подразделения всех передаваемых продаж & HTTPS & 443 \\
Система корпоративного доступа & Выполняет двухфакторную аутентификацию (2FA), проверяет пользовательскую сессию, факт второго фактора и роль. Вход выполняется существующей корпоративной страницей. При недоступности проверки запрос отклоняется & HTTPS & 443 \\
Платформа PostgreSQL & Позволяет создать отдельную логическую БД для нового приложения. Загрузчик читает и записывает, сервис отчётности только читает. Поддерживает заданный объём, ежедневные резервные копии и восстановление в рамках N-02 & PostgreSQL поверх TLS & 5432 \\
Система передачи отчётов & Принимает CSV и возвращает подтверждение внешнего приёма. Имеет обычный канал для общедоступных данных и защищённый канал для конфиденциальных данных: клиентский сертификат и проверка разрешённого получателя. Оба канала работают по HTTPS & HTTPS & 443 \\
Внешняя система-получатель & Принимает отчёт только от корпоративной системы передачи по согласованному защищённому интерфейсу & HTTPS & 443 \\
Технологический прокси & Обязательный выход из контура для внутренних приложений с рабочими часами; HTTPS (HTTP поверх TLS) до прокси на порту 443; метод HTTP CONNECT создаёт туннель к получателю. Внутри проходит отдельный сквозной TLS с клиентским сертификатом системы передачи; прокси не расшифровывает полезную нагрузку & HTTPS & 443 \\
Бизнес-прокси & Выход для приложений, обслуживающих клиентов 24/7; два независимых узла, автоматическое переключение, доступность 99,99\%; поддерживает тот же защищённый туннель & HTTPS & 443 \\
Платформа приложений & Запускает два независимых сервиса, хранит секреты и выполняет задачу по расписанию в загрузчике & HTTPS & 443 \\
\end{longtable}

У издателя есть право записи в поток продаж, у загрузчика --- чтения. Источник справочника доступен только на чтение. Подтверждать обработку сообщения можно после сохранения продажи в БД; номер продажи позволяет безопасно повторить обработку. Учётные данные выдаются платформой и в документах отсутствуют.

\KeepTableStart{1}{%
\begin{tabular}{@{}
  >{\raggedright\arraybackslash}p{(\columnwidth - 2\tabcolsep) * \real{0.2500}}
  >{\raggedright\arraybackslash}p{(\columnwidth - 2\tabcolsep) * \real{0.7500}}@{}}
\toprule\noalign{}
\begin{minipage}[b]{\linewidth}\raggedright
ID
\end{minipage} & \begin{minipage}[b]{\linewidth}\raggedright
Правило
\end{minipage} \\
\midrule\noalign{}
SEC-01 & Все соединения защищены TLS; протоколы и порты берутся из каталога выше \\
\end{tabular}%
}
\subsection{Правила контекста}\label{ux43fux440ux430ux432ux438ux43bux430-ux43aux43eux43dux442ux435ux43aux441ux442ux430}

\begin{longtable}[]{@{}
  >{\raggedright\arraybackslash}p{(\columnwidth - 2\tabcolsep) * \real{0.2500}}
  >{\raggedright\arraybackslash}p{(\columnwidth - 2\tabcolsep) * \real{0.7500}}@{}}
\toprule\noalign{}
\begin{minipage}[b]{\linewidth}\raggedright
ID
\end{minipage} & \begin{minipage}[b]{\linewidth}\raggedright
Правило
\end{minipage} \\
\midrule\noalign{}
\endhead
\bottomrule\noalign{}
\endlastfoot
SEC-01 & Все соединения защищены TLS; протоколы и порты берутся из каталога выше \\
SEC-02 & Сервис формирования отчётности проверяет доступ для каждого пользовательского запроса; браузер не обращается к источникам или БД напрямую \\
SEC-03 & Права сервисов ограничены их назначением; пользовательские сессии и содержимое отчётов не записываются в журналы \\
SEC-04 & Сводка разрешена к передаче только назначенной внешней системе. Для её класса выбирается защищённый канал с клиентским сертификатом и проверкой получателя; обычный канал не допускается \\
SEC-05 & Доступ пользователя к персональным, клиентским и иным данным требует 2FA. Для данного примера это относится и к финансовой сводке без персональных данных; сессия без подтверждённого второго фактора отклоняется \\
SEC-06 & Любой запрос и передача данных за пределы контура проходят через корпоративный прокси; прямой выход запрещён. Это относится также к ответам и служебным данным внешнего обмена \\
SEC-07 & Для приложений, обслуживающих клиентов 24/7, обязателен отказоустойчивый бизнес-прокси; для остальных --- технологический. Профиль приложения определяется требованиями, а не наличием доступного прокси \\
SEC-08 & Персональные, клиентские и критические данные при передаче за пределы контура шифруются либо передаются по защищённому каналу IPsec. В данном контексте конфиденциальные финансовые данные также требуют этой защиты; HTTPS с TLS и сквозным клиентским сертификатом удовлетворяет требованию шифрования \\
ARCH-01 & Используются готовые корпоративные сервисы, брокер и платформа БД \\
ARCH-02 & Дневная загрузка отделяется от обслуживания пользовательских запросов \\
VIEW-01 & Системы и сервисы в логическом представлении --- Application Component; хранимые данные --- Data Object; передача --- Flow, чтение и запись --- Access. В техническом представлении ПО СУБД --- System Software; направленная Association показывает инициатора обращения \\
\end{longtable}

Исключений нет. В этом примере Kafka выбрана как доступный канал передачи продаж. Загрузка по расписанию сама по себе не требует брокера сообщений.

Признак клиентского обслуживания 24/7 отсутствует: N-02 задаёт внутренний пользовательский сервис с рабочими часами, поэтому применяется технологический прокси. Бизнес-прокси включён в каталог как доступный альтернативный актив, но не включается в рабочие маршруты и диаграммы данного среза. Маршрут внешней доставки: система передачи отчётов → технологический прокси → внешний получатель; подтверждение возвращается тем же защищённым путём. Туннель не меняет логического отправителя: внешний получатель принимает отчёт от системы передачи по её клиентскому сертификату. При недоступности требуемых 2FA, прокси или защиты обход запрещён.

Все правила контекста обязательны; исключений не предусмотрено. Неполные вводные и отрицательные сценарии проверяются отдельно от этого полного примера.

\section{Ожидаемый архитектурный результат}\label{ux43eux436ux438ux434ux430ux435ux43cux44bux439-ux440ux435ux437ux443ux43bux44cux442ux430ux442-ux434ux43bux44f-ux441ux438ux441ux442ux435ux43cux44b-ux43eux442ux447ux451ux442ux43dux43eux441ux442ux438}

Решения предложены для ревью; согласование и ответы человека не имитируются.

\subsection{Реестр требований и уточнений --- A-01}\label{ux440ux435ux435ux441ux442ux440-ux442ux440ux435ux431ux43eux432ux430ux43dux438ux439-ux438-ux443ux442ux43eux447ux43dux435ux43dux438ux439--a-01}

Последовательность обработки соответствует FR-01--FR-05 основного BRD. Проектируются два сервиса: загрузка данных и формирование отчётности. Внешние системы и платформы переиспользуются; логическая БД отчётности создаётся на существующей платформе PostgreSQL. Пользователь --- сотрудник с ролью «Аналитик».

\KeepTableStart{0}{%
\begin{tabular}{@{}
  >{\raggedright\arraybackslash}p{(\columnwidth - 4\tabcolsep) * \real{0.3100}}
  >{\raggedright\arraybackslash}p{(\columnwidth - 4\tabcolsep) * \real{0.3500}}
  >{\raggedright\arraybackslash}p{(\columnwidth - 4\tabcolsep) * \real{0.3400}}@{}}
\toprule\noalign{}
\begin{minipage}[b]{\linewidth}\raggedright
Требование
\end{minipage} & \begin{minipage}[b]{\linewidth}\raggedright
Подтверждающий фрагмент BRD
\end{minipage} & \begin{minipage}[b]{\linewidth}\raggedright
Где учтено
\end{minipage} \\
\midrule\noalign{}
R-01 & «Получать новые продажи в течение дня» & Поток через брокер в загрузчик \\
\end{tabular}%
}
\begin{longtable}[]{@{}
  >{\raggedright\arraybackslash}p{(\columnwidth - 4\tabcolsep) * \real{0.3100}}
  >{\raggedright\arraybackslash}p{(\columnwidth - 4\tabcolsep) * \real{0.3500}}
  >{\raggedright\arraybackslash}p{(\columnwidth - 4\tabcolsep) * \real{0.3400}}@{}}
\toprule\noalign{}
\begin{minipage}[b]{\linewidth}\raggedright
Требование
\end{minipage} & \begin{minipage}[b]{\linewidth}\raggedright
Подтверждающий фрагмент BRD
\end{minipage} & \begin{minipage}[b]{\linewidth}\raggedright
Где учтено
\end{minipage} \\
\midrule\noalign{}
\endhead
\bottomrule\noalign{}
\endlastfoot
R-01 & «Получать новые продажи в течение дня» & Поток через брокер в загрузчик \\
R-02 & «Получать полный справочник подразделений каждый день в 06:00» & Обращение загрузчика к системе учёта подразделений \\
R-03 & «Повтор сообщения с тем же номером продажи не создаёт дубль» & Проверка номера и сохранение в БД \\
R-04 & «видит подразделения с продажами» & Чтение данных и подсчёт итогов \\
R-05 & «Файл рассчитывается на момент запроса» & Формирование CSV сервисом отчётности \\
R-06 & «состояние загрузки и время последней успешной обработки» & Служебные сведения в БД и пользовательском ответе \\
R-07 & «назначенной внешней системе» & Защищённый канал через систему передачи отчётов \\
N-01 & «не более 5 секунд» & Целевой показатель; требует испытаний \\
N-02 & «Целевое восстановление --- до 4 часов» & Платформа БД и будущие эксплуатационные проверки \\
N-03 & «Каждый пользовательский запрос разрешён» & Обращение к системе корпоративного доступа, проверка роли и второго фактора по SEC-05 \\
\end{longtable}

Вводные достаточны для проектирования примера; открытых вопросов и принятых допущений нет. Согласование архитектором не моделируется.

\subsection{Каталог данных и матрица потоков --- A-02}\label{ux43aux430ux442ux430ux43bux43eux433-ux434ux430ux43dux43dux44bux445-ux438-ux43cux430ux442ux440ux438ux446ux430-ux43fux43eux442ux43eux43aux43eux432--a-02}

Диаграмма потоков данных --- A-03.

Продажи и сводка --- конфиденциальные финансовые данные без персональных сведений. Справочник, состояние загрузки и результат проверки доступа --- внутренние служебные данные.

\KeepTableStart{0}{%
\begin{tabular}{@{}
  >{\raggedright\arraybackslash}p{(\columnwidth - 4\tabcolsep) * \real{0.3100}}
  >{\raggedright\arraybackslash}p{(\columnwidth - 4\tabcolsep) * \real{0.3500}}
  >{\raggedright\arraybackslash}p{(\columnwidth - 4\tabcolsep) * \real{0.3400}}@{}}
\toprule\noalign{}
\begin{minipage}[b]{\linewidth}\raggedright
Набор
\end{minipage} & \begin{minipage}[b]{\linewidth}\raggedright
Кто предоставляет или формирует
\end{minipage} & \begin{minipage}[b]{\linewidth}\raggedright
Хранение
\end{minipage} \\
\midrule\noalign{}
Данные о продажах & Система продаж & Через брокер поступают в БД отчётности \\
\end{tabular}%
}
\begin{longtable}[]{@{}
  >{\raggedright\arraybackslash}p{(\columnwidth - 4\tabcolsep) * \real{0.3100}}
  >{\raggedright\arraybackslash}p{(\columnwidth - 4\tabcolsep) * \real{0.3500}}
  >{\raggedright\arraybackslash}p{(\columnwidth - 4\tabcolsep) * \real{0.3400}}@{}}
\toprule\noalign{}
\begin{minipage}[b]{\linewidth}\raggedright
Набор
\end{minipage} & \begin{minipage}[b]{\linewidth}\raggedright
Кто предоставляет или формирует
\end{minipage} & \begin{minipage}[b]{\linewidth}\raggedright
Хранение
\end{minipage} \\
\midrule\noalign{}
\endhead
\bottomrule\noalign{}
\endlastfoot
Данные о продажах & Система продаж & Через брокер поступают в БД отчётности \\
Справочник подразделений & Система учёта подразделений & Копия в БД отчётности \\
Данные отчётности & Сервис загрузки данных & Продажи за 30 дней, справочник и состояние загрузки в БД \\
Сводный отчёт & Сервис формирования отчётности & Не сохраняется; выдаётся сотруднику \\
Состояние загрузки & Сервис загрузки данных & Служебные сведения в БД отчётности \\
Подтверждение передачи & Внешняя система через систему передачи отчётов & Не сохраняется; возвращается аналитику \\
Результат проверки доступа & Система корпоративного доступа: роль и подтверждение второго фактора & Не сохраняется \\
Контекст двухфакторного входа & Браузер сотрудника & Не сохраняется в модели; секреты не включаются \\
Подтверждение двухфакторной сессии & Система корпоративного доступа & Не сохраняется в БД отчётности \\
\end{longtable}

\KeepTableStart{0}{%
\begin{tabular}{@{}
  >{\raggedright\arraybackslash}p{(\columnwidth - 6\tabcolsep) * \real{0.2400}}
  >{\raggedright\arraybackslash}p{(\columnwidth - 6\tabcolsep) * \real{0.2400}}
  >{\raggedright\arraybackslash}p{(\columnwidth - 6\tabcolsep) * \real{0.3000}}
  >{\raggedright\arraybackslash}p{(\columnwidth - 6\tabcolsep) * \real{0.2200}}@{}}
\toprule\noalign{}
\begin{minipage}[b]{\linewidth}\raggedright
Откуда
\end{minipage} & \begin{minipage}[b]{\linewidth}\raggedright
Куда
\end{minipage} & \begin{minipage}[b]{\linewidth}\raggedright
Какие данные
\end{minipage} & \begin{minipage}[b]{\linewidth}\raggedright
Когда
\end{minipage} \\
\midrule\noalign{}
Система продаж & Корпоративный брокер сообщений & Новая продажа & После завершения продажи \\
\end{tabular}%
}
\begin{longtable}[]{@{}
  >{\raggedright\arraybackslash}p{(\columnwidth - 6\tabcolsep) * \real{0.2400}}
  >{\raggedright\arraybackslash}p{(\columnwidth - 6\tabcolsep) * \real{0.2400}}
  >{\raggedright\arraybackslash}p{(\columnwidth - 6\tabcolsep) * \real{0.3000}}
  >{\raggedright\arraybackslash}p{(\columnwidth - 6\tabcolsep) * \real{0.2200}}@{}}
\toprule\noalign{}
\begin{minipage}[b]{\linewidth}\raggedright
Откуда
\end{minipage} & \begin{minipage}[b]{\linewidth}\raggedright
Куда
\end{minipage} & \begin{minipage}[b]{\linewidth}\raggedright
Какие данные
\end{minipage} & \begin{minipage}[b]{\linewidth}\raggedright
Когда
\end{minipage} \\
\midrule\noalign{}
\endhead
\bottomrule\noalign{}
\endlastfoot
Система продаж & Корпоративный брокер сообщений & Новая продажа & После завершения продажи \\
Корпоративный брокер сообщений & Сервис загрузки данных & Данные о продажах & В течение дня \\
Система учёта подразделений & Сервис загрузки данных & Справочник подразделений & При запуске и ежедневно в 06:00 \\
Сервис загрузки данных & Данные отчётности & Проверенные продажи, справочник и состояние загрузки & После обработки входных данных \\
Данные отчётности & Сервис формирования отчётности & Продажи, подразделения и состояние загрузки & При пользовательском запросе \\
Система корпоративного доступа & Сервис формирования отчётности & Результат проверки прав и второго фактора & При каждом пользовательском запросе \\
Сервис формирования отчётности & Браузер сотрудника & Сводка, CSV, состояние загрузки и результат передачи & После проверки доступа \\
Сервис формирования отчётности & Система передачи отчётов & CSV с конфиденциальной сводкой & По команде аналитика \\
Система передачи отчётов & Технологический прокси & CSV в сквозном TLS-туннеле & После проверки получателя \\
Технологический прокси & Внешняя система-получатель & CSV в том же защищённом туннеле & При внешней доставке \\
Внешняя система-получатель & Технологический прокси & Подтверждение передачи & Ответ по тому же защищённому обращению \\
Технологический прокси & Система передачи отчётов & Подтверждение передачи & Ответ по тому же защищённому обращению \\
Система передачи отчётов & Сервис формирования отчётности & Подтверждение передачи & После ответа получателя \\
Браузер сотрудника & Система корпоративного доступа & Контекст двухфакторного входа & На существующей странице входа \\
Система корпоративного доступа & Браузер сотрудника & Подтверждение двухфакторной сессии & После успешного второго фактора \\
\end{longtable}

\begin{landscape}\thispagestyle{plain}
\subsection*{Диаграмма потоков данных}
\vfill\begin{center}
\includegraphics[width=\linewidth,height=.82\textheight,keepaspectratio]{.build/documents/data-flow.pdf}
\end{center}\vfill
\end{landscape}

Загрузчик проверяет продажи и сохраняет их без повторов, удаляет устаревшие записи по R-03 и после восстановления. Сервис отчётности ограничивает даты по R-04; запросы за пределами окна отклоняются. Сервис отчётности связывает продажи с наименованиями подразделений, группирует их и подсчитывает количество и сумму. Эти преобразования и чтение/запись БД показаны на уровне назначения сервисов; внутренние алгоритмы и таблицы БД не раскрываются.

\subsection{Каталог систем и матрица интеграций --- A-04}\label{ux43aux430ux442ux430ux43bux43eux433-ux441ux438ux441ux442ux435ux43c-ux438-ux43cux430ux442ux440ux438ux446ux430-ux438ux43dux442ux435ux433ux440ux430ux446ux438ux439--a-04}

Диаграмма межсистемных взаимодействий --- A-05.

Каталог систем задан в YAML: технологический прокси выбран по N-02 и SEC-07; бизнес-прокси остаётся неиспользуемым корпоративным активом. Оба существуют и не разрабатываются в проекте. Создаются только сервис загрузки данных и сервис формирования отчётности. Логическая база отчётности также создаётся. Платформа PostgreSQL и остальные корпоративные системы уже существуют.

\KeepTableStart{0}{%
\begin{tabular}{@{}
  >{\raggedright\arraybackslash}p{(\columnwidth - 8\tabcolsep) * \real{0.2350}}
  >{\raggedright\arraybackslash}p{(\columnwidth - 8\tabcolsep) * \real{0.2350}}
  >{\raggedright\arraybackslash}p{(\columnwidth - 8\tabcolsep) * \real{0.2800}}
  >{\raggedright\arraybackslash}p{(\columnwidth - 8\tabcolsep) * \real{0.1700}}
  >{\raggedright\arraybackslash}p{(\columnwidth - 8\tabcolsep) * \real{0.0800}}@{}}
\toprule\noalign{}
\begin{minipage}[b]{\linewidth}\raggedright
Инициатор
\end{minipage} & \begin{minipage}[b]{\linewidth}\raggedright
Получатель
\end{minipage} & \begin{minipage}[b]{\linewidth}\raggedright
Назначение обращения
\end{minipage} & \begin{minipage}[b]{\linewidth}\raggedright
Протокол
\end{minipage} & \begin{minipage}[b]{\linewidth}\raggedright
Порт
\end{minipage} \\
\midrule\noalign{}
Система продаж & Корпоративный брокер сообщений & Публикация новых продаж & Kafka поверх TLS & 9093 \\
\end{tabular}%
}
\begin{longtable}[]{@{}
  >{\raggedright\arraybackslash}p{(\columnwidth - 8\tabcolsep) * \real{0.2350}}
  >{\raggedright\arraybackslash}p{(\columnwidth - 8\tabcolsep) * \real{0.2350}}
  >{\raggedright\arraybackslash}p{(\columnwidth - 8\tabcolsep) * \real{0.2800}}
  >{\raggedright\arraybackslash}p{(\columnwidth - 8\tabcolsep) * \real{0.1700}}
  >{\raggedright\arraybackslash}p{(\columnwidth - 8\tabcolsep) * \real{0.0800}}@{}}
\toprule\noalign{}
\begin{minipage}[b]{\linewidth}\raggedright
Инициатор
\end{minipage} & \begin{minipage}[b]{\linewidth}\raggedright
Получатель
\end{minipage} & \begin{minipage}[b]{\linewidth}\raggedright
Назначение обращения
\end{minipage} & \begin{minipage}[b]{\linewidth}\raggedright
Протокол
\end{minipage} & \begin{minipage}[b]{\linewidth}\raggedright
Порт
\end{minipage} \\
\midrule\noalign{}
\endhead
\bottomrule\noalign{}
\endlastfoot
Система продаж & Корпоративный брокер сообщений & Публикация новых продаж & Kafka поверх TLS & 9093 \\
Сервис загрузки данных & Корпоративный брокер сообщений & Чтение продаж & Kafka поверх TLS & 9093 \\
Сервис загрузки данных & Система учёта подразделений & Получение справочника & HTTPS & 443 \\
Сервис загрузки данных & База данных отчётности & Сохранение данных и состояния загрузки & PostgreSQL поверх TLS & 5432 \\
Сервис формирования отчётности & База данных отчётности & Чтение данных для сводки и состояния загрузки & PostgreSQL поверх TLS & 5432 \\
Сервис формирования отчётности & Система корпоративного доступа & Проверка роли и второго фактора & HTTPS & 443 \\
Браузер сотрудника & Сервис формирования отчётности & Просмотр, скачивание и передача отчёта & HTTPS & 443 \\
Сервис формирования отчётности & Система передачи отчётов & Защищённая передача CSV и получение подтверждения & HTTPS & 443 \\
Система передачи отчётов & Технологический прокси & Метод CONNECT создаёт туннель к получателю; отчёт и подтверждение внутри сквозного TLS & HTTPS (HTTP поверх TLS) & 443 \\
Технологический прокси & Внешняя система-получатель & Пересылка сквозного TLS-туннеля и возврат подтверждения & HTTPS & 443 \\
Браузер сотрудника & Система корпоративного доступа & Двухфакторный вход & HTTPS & 443 \\
\end{longtable}

\clearpage
\subsection*{Диаграмма межсистемных взаимодействий}
\vfill\begin{center}
\includegraphics[width=\linewidth,height=.82\textheight,keepaspectratio]{.build/documents/interactions.pdf}
\end{center}\vfill
\clearpage

Из двух каналов системы передачи выбирается защищённый: класс сводки требует клиентского сертификата и проверки разрешённого получателя (SEC-04). Прямого соединения системы передачи с внешним получателем нет: весь обмен идёт через технологический прокси (SEC-06, SEC-07). Соединение системы передачи с прокси защищено HTTPS (HTTP поверх TLS); метод HTTP CONNECT создаёт прокси-туннель к разрешённому получателю. Внутри туннеля устанавливается отдельное сквозное HTTPS/TLS-соединение системы передачи с получателем, с клиентским сертификатом системы передачи. CONNECT не является протоколом шифрования: защиту CSV обеспечивает сквозной TLS, который прокси не расшифровывает (SEC-08). 2FA выполняется системой корпоративного доступа, а признак второго фактора и роль проверяются для каждого пользовательского запроса (SEC-02, SEC-05). Подтверждение внешнего приёма возвращается по тому же туннелю и обращению; при его отсутствии интерфейс показывает «Передача не подтверждена».

Загрузчик инициирует чтение из брокера, хотя сами данные идут из брокера в загрузчик. Аналогично, запрос справочника направлен к системе учёта подразделений, а справочник возвращается из неё.

\subsection{Архитектурный манифест --- A-06}\label{ux430ux440ux445ux438ux442ux435ux43aux442ux443ux440ux43dux44bux439-ux43cux430ux43dux438ux444ux435ux441ux442--a-06}

{архитектурный манифест (полностью приведён в приложении)} содержит контекст, понятные названия и назначения систем, описание данных, потоки и технические обращения. Ключи связей осмысленные; служебные номера на диаграммах не показываются. Реестр требований и уточнений (A-01) и Реестр решений и отчёт проверок (A-07) остаются в этом документе. Справочник \texttt{источники} указывает идентификатор документа, путь, раздел и редакцию (полный hash содержащего комплект коммита); происхождение каждого объекта и связи содержит ссылки на справочник, конкретные требования/правила или строку каталога и автора «Ручной эталон». Повторяющиеся свойства вынесены в именованные профили; у каждого объекта и связи остаются явная ссылка на профиль происхождения и собственные основания. Это не имитация вывода ИИ.

Объекты и связи имеют объяснимые статусы контроля по FR-09. Для объектов и связей задан базовый статус «Требует проверки»: основания перечислены, но трактовка и решения ещё не подтверждены архитектором. Если обязательного основания не хватает, применяется «Недостаточно данных»; доказанный дефект получает статус «Нарушение». Статус «Подтверждено» требует актуальных проверенных оснований и закрытия обязательных условий; сам по себе он не означает согласования. При изменении источника зависимые подтверждения проверяются повторно.

Для A-03 (диаграмма потоков данных) и A-05 (диаграмма межсистемных взаимодействий) используются ArchiMate и профиль VIEW-01. Обязательные исходники: {data-flow.puml} и {interactions.puml}; SVG сформированы из них. Бизнес-прокси и платформа приложений не входят в срез действующих потоков; платформа PostgreSQL показана группировкой БД в техническом представлении. Исключение из среза не удаляет актив из корпоративного каталога.

\KeepTableStart{2}{%
\begin{tabular}{@{}
  >{\raggedright\arraybackslash}p{(\columnwidth - 4\tabcolsep) * \real{0.3100}}
  >{\raggedright\arraybackslash}p{(\columnwidth - 4\tabcolsep) * \real{0.3500}}
  >{\raggedright\arraybackslash}p{(\columnwidth - 4\tabcolsep) * \real{0.3400}}@{}}
\toprule\noalign{}
\begin{minipage}[b]{\linewidth}\raggedright
Предложение и причина
\end{minipage} & \begin{minipage}[b]{\linewidth}\raggedright
Существенная альтернатива
\end{minipage} & \begin{minipage}[b]{\linewidth}\raggedright
Авторитетный источник
\end{minipage} \\
\midrule\noalign{}
Отделить загрузчик от пользовательского сервиса: продажи поступают в рабочее время & Совместить обработку и пользовательские запросы; противоречит правилу разделения & R-01, ARCH-02 \\
\end{tabular}%
}
\subsection{Реестр решений и отчёт проверок --- A-07}\label{ux440ux435ux435ux441ux442ux440-ux440ux435ux448ux435ux43dux438ux439-ux438-ux43eux442ux447ux451ux442-ux43fux440ux43eux432ux435ux440ux43eux43a--a-07}

\subsubsection{Решения}\label{ux440ux435ux448ux435ux43dux438ux44f}

\begin{longtable}[]{@{}
  >{\raggedright\arraybackslash}p{(\columnwidth - 4\tabcolsep) * \real{0.3100}}
  >{\raggedright\arraybackslash}p{(\columnwidth - 4\tabcolsep) * \real{0.3500}}
  >{\raggedright\arraybackslash}p{(\columnwidth - 4\tabcolsep) * \real{0.3400}}@{}}
\toprule\noalign{}
\begin{minipage}[b]{\linewidth}\raggedright
Предложение и причина
\end{minipage} & \begin{minipage}[b]{\linewidth}\raggedright
Существенная альтернатива
\end{minipage} & \begin{minipage}[b]{\linewidth}\raggedright
Авторитетный источник
\end{minipage} \\
\midrule\noalign{}
\endhead
\bottomrule\noalign{}
\endlastfoot
Отделить загрузчик от пользовательского сервиса: продажи поступают в рабочее время & Совместить обработку и пользовательские запросы; противоречит правилу разделения & R-01, ARCH-02 \\
Получать продажи через готовый Kafka, а справочник --- по расписанию & Периодически запрашивать продажи по HTTP; такого интерфейса во вводных нет & R-01, R-02; каталог систем \\
Сохранять данные в готовой PostgreSQL и учитывать номер продажи перед подтверждением сообщения & Формировать отчёт прямо из источников; не даёт собственного накопленного набора и защиты от повторов & R-03, R-06, ARCH-01 \\
Использовать защищённый канал для передачи сводки внешнему получателю & Обычный канал предназначен только для общедоступных данных & R-07, N-03, SEC-04 \\
Проверять роль и второй фактор через готовую систему корпоративного доступа & Собственный механизм входа; оснований для новой системы нет & N-03, SEC-02, SEC-05 \\
Использовать технологический прокси и сквозной TLS для внешнего обмена & Прямой выход запрещён; бизнес-прокси предназначен для клиентского обслуживания 24/7 & N-02, SEC-06, SEC-07, SEC-08 \\
\end{longtable}

\KeepTableStart{1}{%
\begin{tabular}{@{}
  >{\raggedright\arraybackslash}p{(\columnwidth - 4\tabcolsep) * \real{0.3100}}
  >{\raggedright\arraybackslash}p{(\columnwidth - 4\tabcolsep) * \real{0.3500}}
  >{\raggedright\arraybackslash}p{(\columnwidth - 4\tabcolsep) * \real{0.3400}}@{}}
\toprule\noalign{}
\begin{minipage}[b]{\linewidth}\raggedright
Что проверено
\end{minipage} & \begin{minipage}[b]{\linewidth}\raggedright
По какому источнику
\end{minipage} & \begin{minipage}[b]{\linewidth}\raggedright
Результат
\end{minipage} \\
\midrule\noalign{}
Синтаксис YAML, структура эталона, типы полей и ссылок, названия, таблицы и исходники диаграмм & Структура ручного эталона; FR-06 & Проверено при подготовке комплекта; преобразование в Python-доменные классы и валидаторы модели предстоит реализовать \\
\end{tabular}%
}
\subsubsection{Проверки}\label{ux43fux440ux43eux432ux435ux440ux43aux438}

\begin{longtable}[]{@{}
  >{\raggedright\arraybackslash}p{(\columnwidth - 4\tabcolsep) * \real{0.3100}}
  >{\raggedright\arraybackslash}p{(\columnwidth - 4\tabcolsep) * \real{0.3500}}
  >{\raggedright\arraybackslash}p{(\columnwidth - 4\tabcolsep) * \real{0.3400}}@{}}
\toprule\noalign{}
\begin{minipage}[b]{\linewidth}\raggedright
Что проверено
\end{minipage} & \begin{minipage}[b]{\linewidth}\raggedright
По какому источнику
\end{minipage} & \begin{minipage}[b]{\linewidth}\raggedright
Результат
\end{minipage} \\
\midrule\noalign{}
\endhead
\bottomrule\noalign{}
\endlastfoot
Синтаксис YAML, структура эталона, типы полей и ссылок, названия, таблицы и исходники диаграмм & Структура ручного эталона; FR-06 & Проверено при подготовке комплекта; преобразование в Python-доменные классы и валидаторы модели предстоит реализовать \\
Направления данных и технических обращений разделены; загрузчик выделен & R-01, ARCH-02, VIEW-01 & Пройдено \\
Протоколы защищены; порты соответствуют заданному каталогу & SEC-01 & Пройдено \\
Права доступа, защита от дублей и сохранность сообщений предусмотрены & R-03, R-06, SEC-02, SEC-03 & Учтено в проекте; выполнение требует тестов реализации \\
2FA и проверка второго фактора в сессии & SEC-02, SEC-05 & Учтено в модели; реальная аутентификация не испытывалась \\
Внешний получатель, прокси, классификация и защищённый туннель & R-07, N-02, N-03, SEC-04, SEC-06--SEC-08 & Учтено в проекте; доставка, доступность и сертификаты требуют испытаний \\
Происхождение объектов и связей, статусы контроля & FR-09 & Основания указаны; статус «Требует проверки», решения ожидают архитектора; метрики генерации не измерялись \\
Граф данных и вызовов, законные граничные роли & FR-07, VIEW-01 & Пройдено для выбранного среза с учётом ответов, чтения/записи и границ; система продаж и справочник --- заданные источники, БД --- хранилище, получатель --- внешний участник \\
Типы элементов и связей соответствуют выбранному профилю & VIEW-01 & Пройдено для этого профиля \\
Задержка данных, время ответа, доступность и восстановление достигнуты & R-01, N-01, N-02 & Не проверено: нужны эксплуатационные испытания \\
\end{longtable}

Полная нормативная проверка ArchiMate и согласование человеком остаются открытыми. Структурная проверка эталона не доказывает работоспособность будущей реализации.

\section{Полный архитектурный манифест}
\label{manifest-appendix}
Ниже приведён полный YAML-манифест системы формирования отчётности.
Он описывает контекст, системы, наборы данных, потоки и технические взаимодействия.
\VerbatimInput[fontsize=\normalsize,breaklines=true,breakanywhere=true,
  breaksymbolleft={},breaksymbolright={},tabsize=2,
  frame=leftline,framerule=.6pt,rulecolor=\color{Accent},
  xleftmargin=8pt,framesep=8pt]{docs/examples/reporting-system/manifest.yaml}

\end{document}
